\documentclass[../../main.tex]{subfiles}

\begin{document}

\chapter{Kriptografi Modern Simetris: Stream Cipher}

\begin{learningoutcome}
Setelah mempelajari bab ini, mahasiswa diharapkan mampu:
\begin{itemize}
    \item Menjelaskan perbedaan mendasar antara kriptografi klasik dan modern, terutama dalam representasi data dan operasi dasar (Sub-CPMK 1.2).
    \item Memahami konsep operasi bitwise XOR dan aplikasinya dalam algoritma kriptografi (Sub-CPMK 1.2).
    \item Menjelaskan mekanisme kerja \textit{Stream Cipher} dan peran \textit{keystream generator} (Sub-CPMK 1.2).
    \item Menganalisis algoritma \textit{stream cipher} populer seperti RC4, A5/1, Trivium, dan Salsa20/ChaCha20 (Sub-CPMK 1.2).
    \item Membedakan antara \textit{Stream Cipher} dan \textit{Block Cipher} berdasarkan cara pemrosesan data (Sub-CPMK 1.2).
\end{itemize}
\end{learningoutcome}

\subfile{section-01}
\subfile{section-02}
\subfile{section-03}
\section{Contoh Terhitung}

\begin{example}
\textbf{Enkripsi dengan \textit{stream cipher}.} Plainteks dan keystream masing-masing 24 bit, ditulis dalam heksadesimal sebagai $P=\code{571964}$ dan $K=\code{F5923D}$. Nyatakan keduanya dalam biner, lalu XOR bit per bit, $c_i=p_i\oplus k_i$:
\[
\begin{array}{rcl}
P &=& 0101\,0111\,0001\,1001\,0110\,0100\\
K &=& 1111\,0101\,1001\,0010\,0011\,1101\\
\hline
C &=& 1010\,0010\,1000\,1011\,0101\,1001
\end{array}
\]
sehingga cipherteksnya adalah $C=\code{A28B59}$. Dekripsi memakai keystream yang sama:
\[
C\oplus K=\code{A28B59}\oplus\code{F5923D}=\code{571964}=P.
\]
Karena tiap bit diproses sendiri-sendiri, galat satu bit pada cipherteks hanya merusak satu bit plainteks hasil dekripsi: mengubah bit ke-5 pada $C$ menjadi $\code{AA8B59}$ memberi $P'=\code{5F1964}$, yang berbeda dari $P$ tepat pada bit ke-5 saja.
\end{example}

\begin{example}
\textbf{Membangkitkan keystream dengan LFSR.} Pakailah LFSR 4 bit pada
Bagian~\ref{sec:konsep-stream-cipher}, yaitu umpan balik $f=b_3\oplus b_0$
dengan muatan awal $1111$. Tabel~\ref{tab:lfsr-4bit} menunjukkan bahwa barisan
keluarannya berperiode 15 bit:
\[
\code{1111\,0101\,1001\,000}\ .
\]
Karena pesannya lebih panjang daripada periode itu, keystream diulang dari awal.
Untuk pesan \code{Aku} yang panjangnya 24 bit, keystream yang dipakai adalah
15 bit pertama ditambah 9 bit awalnya lagi:
\[
K=\code{1111\,0101\,1001\,0001\,1110\,1011}=\code{F591EB}.
\]
Enkripsi dilakukan dengan XOR per bit:
\[
\begin{array}{rcl}
P &=& 0100\,0001\,0110\,1011\,0111\,0101 \quad(\code{Aku})\\
K &=& 1111\,0101\,1001\,0001\,1110\,1011\\
\hline
C &=& 1011\,0100\,1111\,1010\,1001\,1110
\end{array}
\]
sehingga cipherteksnya adalah $\code{B4FA9E}$. Dekripsi memakai keystream yang
sama: $\code{B4FA9E}\oplus\code{F591EB}=\code{416B75}$, yaitu \code{Aku} kembali.

Perhatikan pengulangan yang tampak pada baris $K$: sembilan bit pertama
keystream muncul kembali pada sembilan bit terakhir, karena periode keystream
hanya 15 bit sedangkan pesannya 24 bit. Pengulangan inilah yang harus dihindari
oleh keystream generator yang baik, dan inilah alasan pembahasan pada
Bagian~\ref{sec:konsep-stream-cipher} menekankan pentingnya periode yang jauh
lebih panjang daripada pesan.
\end{example}

\begin{example}
\textbf{RC4 versi mainan ($N=8$) untuk kunci \code{AB}.} Agar seluruh langkah
dapat diperiksa dengan tangan, larik $S$ diperkecil dari 256 menjadi 8 elemen,
sehingga semua operasi modulo 256 menjadi modulo 8. Kunci \code{AB} memberi
byte $\code{41}_{16}=65$ dan $\code{42}_{16}=66$, yang setelah direduksi modulo
8 menjadi $K=[\,1,\ 2\,]$; kunci diulang periodik: $K_0=1$, $K_1=2$, $K_2=1$,
dan seterusnya.

\textbf{KSA.} Larik $S$ mula-mula berisi $[\,0,1,2,3,4,5,6,7\,]$. Penunjuk $j$
bergerak menurut $j \leftarrow (j + S_i + K_{i \bmod 2}) \bmod 8$, lalu $S_i$
dan $S_j$ dipertukarkan. Hasil setiap langkah adalah
\[
\begin{array}{cll}
i=0: & j = (0+0+1) \bmod 8 = 1 & \text{swap } S_0 \leftrightarrow S_1
      \rightarrow [\,1,0,2,3,4,5,6,7\,]\\
i=1: & j = (1+0+2) \bmod 8 = 3 & \text{swap } S_1 \leftrightarrow S_3
      \rightarrow [\,1,3,2,0,4,5,6,7\,]\\
i=2: & j = (3+2+1) \bmod 8 = 6 & \text{swap } S_2 \leftrightarrow S_6
      \rightarrow [\,1,3,6,0,4,5,2,7\,]\\
i=3: & j = (6+0+2) \bmod 8 = 0 & \text{swap } S_3 \leftrightarrow S_0
      \rightarrow [\,0,3,6,1,4,5,2,7\,]\\
i=4: & j = (0+4+1) \bmod 8 = 5 & \text{swap } S_4 \leftrightarrow S_5
      \rightarrow [\,0,3,6,1,5,4,2,7\,]\\
i=5: & j = (5+4+2) \bmod 8 = 3 & \text{swap } S_5 \leftrightarrow S_3
      \rightarrow [\,0,3,6,4,5,1,2,7\,]\\
i=6: & j = (3+2+1) \bmod 8 = 6 & \text{swap } S_6 \leftrightarrow S_6
      \rightarrow [\,0,3,6,4,5,1,2,7\,]\\
i=7: & j = (6+7+2) \bmod 8 = 7 & \text{swap } S_7 \leftrightarrow S_7
      \rightarrow [\,0,3,6,4,5,1,2,7\,]
\end{array}
\]
Dua langkah terakhir tidak mengubah apa pun karena indeks yang dipertukarkan
sama. Setelah KSA, larik $S$ menjadi $[\,0,3,6,4,5,1,2,7\,]$ --- inilah
permutasi yang dipakai PRGA.

\textbf{PRGA.} Penunjuk $i$ dan $j$ direset ke nol, lalu untuk setiap byte
keystream dihitung $i \leftarrow (i+1) \bmod 8$, $j \leftarrow (j+S_i) \bmod 8$,
pertukaran $S_i \leftrightarrow S_j$, dan $t \leftarrow (S_i+S_j) \bmod 8$.
Byte keystream adalah $S_t$. Untuk larik $S$ di atas hasilnya adalah
\[
\begin{array}{c|c|c|c|c}
\text{langkah} & i & j & \text{swap} & t \ \rightarrow\ u\\
\hline
1 & 1 & 3 & S_1 \leftrightarrow S_3 & (4+3) \bmod 8 = 7 \rightarrow S_7 = 7\\
2 & 2 & 1 & S_2 \leftrightarrow S_1 & (4+6) \bmod 8 = 2 \rightarrow S_2 = 4\\
3 & 3 & 5 & S_3 \leftrightarrow S_5 & (5+3) \bmod 8 = 0 \rightarrow S_0 = 0\\
4 & 4 & 7 & S_4 \leftrightarrow S_7 & (7+3) \bmod 8 = 2 \rightarrow S_2 = 4\\
5 & 5 & 0 & S_5 \leftrightarrow S_0 & (0+1) \bmod 8 = 1 \rightarrow S_1 = 6
\end{array}
\]
Jadi tiga byte keystream pertama adalah $7$, $4$, dan $0$. Bila plainteksnya
\code{Halo}, yaitu byte $\code{48}$, $\code{61}$, dan $\code{6C}$, maka
\[
\begin{array}{c|ccc}
\text{plainteks} & \code{48} & \code{61} & \code{6C} \quad(\text{\code{Halo}})\\
\text{keystream} & \code{07} & \code{04} & \code{00}\\
\hline
\text{cipherteks} & \code{4F} & \code{65} & \code{6C}
\end{array}
\]
Cipherteks \code{4F656C} tampak seperti data acak, dan dekripsi mengulang
prosedur PRGA yang sama persis untuk mengembalikan \code{Halo}.

Contoh ini sekaligus memperlihatkan mengapa ukuran larik $S$ yang sebenarnya
dipilih 256 dan bukan 8. Dengan $N=8$ hanya ada $8! = 40.320$ permutasi larik
yang mungkin, dan setelah KSA larik $S$ hasilnya dapat dicari dengan menelusuri
seluruh kemungkinan itu. Pada RC4 yang sesungguhnya, $N = 256$ sehingga jumlah
permutasi mencapai $256!$ --- bilangan dengan lebih dari 500 digit --- dan
penelusuran menyeluruh menjadi mustahil. Yang menjatuhkan RC4 bukanlah
kekurangan ukuran lariknya, melainkan korelasi antara byte awal keystream dan
kunci.
\end{example}

\begin{example}
\textbf{Satu \textit{quarter round} ChaCha20.} Kerjakan QR ChaCha20 atas
empat kata $a=1$, $b=2$, $c=3$, dan $d=4$ (semuanya dalam heksadesimal; tiap
kata adalah bilangan 32 bit, sehingga batas penjumlahannya adalah
$2^{32}=\code{FFFFFFFF}$). Urutan operasinya adalah penjumlahan, XOR, dan
rotasi seperti yang tertulis pada Bagian~\ref{sec:algoritma-stream-cipher}.
Delapan langkah yang harus dikerjakan adalah sebagai berikut:
\[
\begin{array}{lr}
a \leftarrow a+b                          & \code{00000003}\\
d \leftarrow \mathrm{rotl}(d \oplus a,16) & \code{00070000}\\
c \leftarrow c+d                          & \code{00070003}\\
b \leftarrow \mathrm{rotl}(b \oplus c,12) & \code{70001000}\\
a \leftarrow a+b                          & \code{70001003}\\
d \leftarrow \mathrm{rotl}(d \oplus a,8)  & \code{07100370}\\
c \leftarrow c+d                          & \code{07170373}\\
b \leftarrow \mathrm{rotl}(b \oplus c,7)  & \code{8B89B9BB}
\end{array}
\]
Sebagai pemeriksaan, kerjakan langkah ketiga dan keempat secara rinci. Pada
saat itu $a=\code{00000003}$ dan $d=\code{00070000}$, sehingga
\[
c = \code{00000003} + \code{00070000} = \code{00070003},
\]
lalu
\[
b = \mathrm{rotl}(\code{00000002} \oplus \code{00070003},\ 12)
  = \mathrm{rotl}(\code{00070001},\ 12)
  = \code{70001000}.
\]
Hasil akhir seluruh QR adalah $a=\code{70001003}$, $b=\code{8B89B9BB}$,
$c=\code{07170373}$, dan $d=\code{07100370}$.

Perhatikan dua hal. Pertama, masukan yang sangat sederhana ($1, 2, 3, 4$)
berubah menjadi empat kata yang tampak acak hanya dalam satu QR. Kedua,
$\mathrm{rotl}(x,n)$ berarti memutar bit $x$ ke kiri sejauh $n$ posisi: bit yang
keluar dari ujung kiri masuk kembali dari ujung kanan, berbeda dari pergeseran
biasa yang membuang bit tersebut. Pada rotasi 16 bit misalnya,
$\code{00000007}$ menjadi $\code{00070000}$ karena tiga bit terendah yang
bernilai $1$ berpindah tempat, bukan hilang.

Contoh ini menunjukkan mengapa ChaCha20 disebut memakai operasi AXR: seluruh
langkah di atas hanya memakai penjumlahan modulo $2^{32}$, XOR, dan rotasi ---
tidak ada perkalian, tidak ada tabel substitusi, dan tidak ada pembagian.
Konsekuensinya, satu QR dapat dikerjakan prosesor dengan beberapa instruksi
saja, dan itulah sumber kecepatan ChaCha20 pada perangkat lunak.
\end{example}


\begin{obeactivity}
\textbf{Aktivitas 5.1: Simulasi XOR Cipher}
1. Tuliskan pesan singkat dalam bentuk biner (8-bit per karakter).
2. Buatlah kunci biner dengan panjang yang sama.
3. Lakukan operasi XOR untuk mendapatkan cipherteks.
4. Berikan cipherteks tersebut kepada rekan Anda dan minta mereka mendekripsinya menggunakan kunci yang sama.
5. Diskusikan apa yang terjadi jika satu bit pada cipherteks diubah (bit-flip).
\end{obeactivity}


\begin{obereflection}
\textbf{Latihan 5.1}
1. Jika plainteks adalah \texttt{10110010} dan kunci adalah \texttt{11001010}, hitunglah cipherteksnya.
2. Mengapa \textit{stream cipher} lebih cocok digunakan untuk komunikasi suara (VoIP) dibandingkan \textit{block cipher}?
3. Jelaskan mengapa penggunaan kunci yang sama untuk dua pesan berbeda dalam \textit{stream cipher} (disebut \textit{key reuse}) sangat berbahaya.
4. Nyatakan kalimat \texttt{Kripto} dalam bentuk biner ASCII, lalu tuliskan pula bentuk heksadesimalnya.
5. Sebuah LFSR 4 bit bermuatan awal \texttt{1010} memakai umpan balik $f=b_3\oplus b_0$. Tentukan isi register pada empat detak pertama beserta bit keluarannya.
6. Mengapa LFSR tunggal tidak aman dipakai sebagai \textit{keystream generator}, padahal periode keluarannya dapat mencapai $2^n-1$ bit?
7. Jelaskan perbedaan peran \textit{key}, \textit{nonce}, dan \textit{counter} pada Salsa20. Manakah yang harus dirahasiakan, dan mengapa?
8. Bandingkan Salsa20 dan ChaCha20 dari segi bentuk \textit{quarter round} serta akibatnya pada penyebaran bit.
\end{obereflection}

\begin{obereflection}
\textbf{Latihan 5.2: Hitungan dan Analisis}
1. Sebuah pesan 12 bit \texttt{011011010110} dienkripsi dengan LFSR 4 bit bermuatan awal \texttt{1111} dan umpan balik $f=b_3\oplus b_0$. Tuliskan keystream 12 bit yang dipakai, lalu tentukan cipherteksnya.
2. Pada RC4 versi mainan dengan $N=8$ dan kunci \texttt{B}, tentukan larik $S$ setelah KSA, lalu hitung byte keystream pertama.
3. Frame GSM berukuran 228 bit dikirim setiap 4,6 milidetik. Hitunglah laju bit yang harus dihasilkan A5/1 untuk mengikuti aliran itu. Mengapa \textit{warm-up} 100 putaran tidak mengganggu laju tersebut?
4. Trivium memakai tiga NLFSR berukuran 93, 84, dan 111 bit dengan total keadaan 288 bit. Jelaskan mengapa penambahan gerbang AND membuat serangan berbasis persamaan linier tidak lagi bekerja.
5. Bandingkan jumlah \textit{warm-up clock} A5/1 dan Trivium relatif terhadap jumlah bit yang dihasilkan per siklus. Apa alasan perbedaan itu?
6. Seorang penyerang menangkap tiga cipherteks TLS yang seluruhnya memakai \textit{nonce} ChaCha20 yang sama tetapi kunci berbeda. Apakah penyerang memperoleh keuntungan? Jelaskan.
7. Seorang penyerang menangkap dua cipherteks $c_1$ dan $c_2$ yang dibangkitkan dengan keystream yang sama, serta berhasil menebak bahwa plainteks pertama $p_1$ dimulai dengan kata \texttt{LAPORAN}. Tunjukkan bagaimana $p_2$ dapat dipulihkan tanpa mengetahui kuncinya.
\end{obereflection}

\section{Rangkuman}
\begin{itemize}
    \item Kriptografi modern beroperasi pada bit/byte, bukan karakter; XOR menjadi operasi dasar karena $a\oplus a=0$ membuat enkripsi dan dekripsi memakai fungsi yang sama.
    \item Satu byte terdiri atas 8 bit dan satu digit heksadesimal mewakili 4 bit, sehingga satu byte dituliskan sebagai dua digit heksadesimal. Penulisan heksadesimal dipakai setiap kali nilai bit perlu disebutkan tanpa menuliskan seluruh rangkaiannya.
    \item \textit{Stream cipher} mengenkripsi satu bit/byte sekaligus dengan $c_i=p_i\oplus k_i$, sehingga cepat dan cocok untuk data \textit{real-time}.
    \item Keamanannya bergantung pada \textit{keystream generator} (PRNG) yang membangkitkan keystream sepanjang pesan dari kunci dan IV sebagai \textit{seed}; penerima harus dapat membangkitkan keystream yang identik.
    \item Keystream tidak boleh dipakai ulang: dua cipherteks berkeystream sama dapat di-XOR sehingga XOR kedua plainteksnya terbuka tanpa mengetahui kunci.
    \item Galat 1 bit pada cipherteks hanya merusak 1 bit plainteks (tidak menyebar), berbeda dengan \textit{block cipher} yang menyebarkan galat antar blok. Karena itu cipher alir tidak menjamin keutuhan data.
    \item LFSR adalah \textit{keystream generator} paling sederhana: tiap detak menggeser isi register satu posisi dan bit umpan balik mengisi posisi yang kosong. LFSR $n$-bit berperiode maksimum $2^n-1$ bila polinomial umpan baliknya primitif, tetapi LFSR tunggal tidak aman karena keluarannya linier dan dapat dipulihkan dengan algoritma Berlekamp--Massey.
    \item RC4 memakai larik $S$ 256 byte yang diacak oleh KSA dan ditarik oleh PRGA. RC4 tidak lagi aman karena byte awal keystream berkorelasi dengan kunci; implementasi yang berhati-hati membuang 256--512 byte awal, dan RFC 7465 melarang pemakaiannya di TLS sejak 2015.
    \item A5/1 memakai tiga LFSR berukuran 19, 22, dan 23 bit dengan kaidah mayoritas pada bit kendalinya; satu frame GSM berukuran 228 bit dan dikirim setiap 4,6 milidetik. A5/1 sudah dipecahkan sejak serangan \textit{known-plaintext} Ross Anderson pada 1994.
    \item Trivium memakai tiga NLFSR berukuran 93, 84, dan 111 bit. Gerbang AND membuat hubungannya non-linier, dan pemanasan $4 \times 288 = 1152$ detak mencampur kunci serta IV sebelum satu bit pun keluar.
    \item Salsa20 dan ChaCha20 memakai operasi ARX pada matriks keadaan $4\times4$ berisi 16 kata 32 bit, sehingga satu pemanggilan menghasilkan blok keystream 512 bit. Karena tiap blok bergantung pada kunci, \textit{nonce}, dan \textit{counter}, blok-blok keystream dapat dihitung secara paralel dan diakses langsung; ChaCha20 bersama Poly1305 menjadi cipher standar TLS 1.3.
    \item Pelajaran dari keempat cipher itu: panjang kunci bukan penentu utama keamanan cipher alir. Yang menentukan adalah ada tidaknya hubungan statistik antara kunci dan keystream yang dapat dilacak penyerang.
\end{itemize}


\begin{obeassessment}
\textbf{Kuis Bab 5}
1. Apa perbedaan utama antara \textit{true random} dan \textit{pseudo-random} dalam konteks \textit{keystream generator}?
2. Mengapa operasi XOR dipilih sebagai operasi utama dalam \textit{stream cipher}?
3. Jelaskan prinsip kerja \textit{majority clocking} pada algoritma A5/1.
4. Hitunglah cipherteks dari plainteks \texttt{11010110} yang di-XOR dengan kunci \texttt{01101100}.
5. Sebutkan tiga syarat yang harus dipenuhi agar keystream generator dapat dipakai berulang kali dengan kunci yang sama untuk pesan berbeda.
6. Jelaskan mengapa RC4 yang menerima kunci sampai 2048 bit tetap dapat dipecahkan, sedangkan panjang kunci justru dianggap sebagai ukuran keamanan pada cipher lain.
7. Apa fungsi pemanasan 1152 detak pada Trivium, dan apa yang terjadi bila tahap itu dihilangkan?
\end{obeassessment}
\begin{competencychecklist}
\begin{tabularx}{\linewidth}{|X|c|}
\hline
Indikator Kompetensi & Tercapai \\
\hline
Saya dapat menjelaskan perbedaan kriptografi klasik dan modern & [ ] \\
\hline
Saya dapat mengonversi data antara biner, byte, dan heksadesimal & [ ] \\
\hline
Saya dapat melakukan operasi bitwise XOR & [ ] \\
\hline
Saya dapat menjelaskan mekanisme kerja \textit{stream cipher} & [ ] \\
\hline
Saya dapat menjelaskan peran \textit{key}, \textit{nonce}, dan \textit{counter} pada keystream generator & [ ] \\
\hline
Saya dapat menghitung keystream LFSR beserta periodenya & [ ] \\
\hline
Saya dapat menjelaskan mengapa LFSR tunggal tidak aman & [ ] \\
\hline
Saya dapat menjelaskan tahap KSA dan PRGA pada RC4 serta kelemahannya & [ ] \\
\hline
Saya dapat menjelaskan kaidah mayoritas dan pola inisialisasi A5/1 & [ ] \\
\hline
Saya dapat menjelaskan struktur NLFSR dan pemanasan pada Trivium & [ ] \\
\hline
Saya dapat menjelaskan operasi AXR dan \textit{quarter round} ChaCha20 & [ ] \\
\hline
Saya dapat membedakan karakteristik RC4, A5/1, dan ChaCha20 & [ ] \\
\hline
Saya dapat menganalisis keuntungan \textit{stream cipher} untuk data \textit{real-time} & [ ] \\
\hline
Saya dapat menunjukkan bahaya pemakaian ulang keystream melalui perhitungan XOR & [ ] \\
\hline
\end{tabularx}
\end{competencychecklist}


\end{document}
